\documentclass[aspectratio=169,11pt]{beamer}

% common_setting.tex
% Shared preamble for the five PSPDE2026 workshop beamer decks (00-04).
% \input this right after \documentclass in each deck.
% Requires xelatex (fontspec/unicode-math need XeTeX or LuaTeX, not pdflatex).

% Fonts
\usepackage{fontspec}
\usepackage{unicode-math}
\setmainfont{Liberation Sans}
\setsansfont{Liberation Sans}
\setmonofont{Liberation Mono}[Scale=0.85]
\setmathfont{Latin Modern Math}

% Theme and colors
\usetheme{moloch}
\molochset{
  block=fill,
  progressbar=frametitle,
  sectionpage=progressbar,
}
\setbeamertemplate{page number in head/foot}[totalframenumber]

% Graphics
\usepackage{graphicx}
\graphicspath{{./figures/}}

% Tables
\usepackage{booktabs}
\usepackage{tabularx}

% Unnumbered captions without the "Figure:" label prefix
\usepackage{caption}
\captionsetup{labelformat=empty}

% Code listings
\usepackage{listings}
\usepackage{xcolor}
\definecolor{codebg}{RGB}{245,245,245}
% moloch's own accent colors, aliased under their original Metropolis names
\colorlet{mAlertFg}{mLightBrown}
\colorlet{mExampleFg}{mLightGreen}
\lstdefinestyle{python}{
  language=Python,
  basicstyle=\ttfamily\scriptsize,
  keywordstyle=\color{mAlertFg}\bfseries,
  commentstyle=\color{mExampleFg}\itshape,
  stringstyle=\color{mExampleFg!70!black},
  backgroundcolor=\color{codebg},
  rulecolor=\color{black!15},
  frame=single,
  breaklines=true,
  showstringspaces=false,
  columns=fullflexible,
  numbers=none
}
\lstdefinestyle{shell}{
  basicstyle=\ttfamily\scriptsize,
  backgroundcolor=\color{codebg},
  rulecolor=\color{black!15},
  frame=single,
  breaklines=true,
  columns=fullflexible
}


% Title information
\title[Modern Workflow in Geology]{Modern Workflow in Geology}
\subtitle{Process Structural and Petrological Data Efficiently}
\author[O.\ Lexa]{Ondrej Lexa}
\institute[Charles University]{
  Institute of Petrology and Structural Geology (IPSG)\\
  Charles University, Prague, Czechia
}
\date{8 October 2026}

\begin{document}

\begin{frame}[plain]
  \titlepage
\end{frame}

\section{Welcome}

\begin{frame}{Today's Goal}
  \begin{block}{Take home ready-made notebooks}
    Jupyter notebooks into which you simply ``pour'' your own Excel tables ---
    structural measurements or microprobe data --- and get stereograms,
    recalculations and charts.
  \end{block}
  \vspace{0.3cm}
  \textbf{Assumed:} basic Python, Jupyter Lab, and how a pandas
  \texttt{DataFrame} works.\\[0.2cm]
  \textbf{Tools:} \texttt{apsg} (structural geology) and \texttt{petropandas}
  (petrology), both built on the scientific Python stack.
\end{frame}

\begin{frame}{Today's Schedule}
  \centering
  \small
  \begin{tabularx}{0.92\textwidth}{lX}
    \toprule
    \textbf{Time} & \textbf{Block} \\
    \midrule
    09:00--09:30 & Introduction and preparation of the working environment \\
    09:30--11:00 & \textbf{Block 1:} Structural geology with \texttt{apsg} \\
    11:00--11:15 & \emph{Coffee break} \\
    11:15--12:30 & \textbf{Block 2:} Advanced structural and tensor analysis \\
    12:30--13:30 & \emph{Lunch break} \\
    13:30--15:00 & \textbf{Block 3:} Processing petrological data with \texttt{petropandas} \\
    15:00--15:15 & \emph{Coffee break} \\
    15:15--16:30 & \textbf{Block 4:} Bring your own data (BYOD) \\
    \bottomrule
  \end{tabularx}
\end{frame}

\section{Scientific Python}

\begin{frame}{The Scientific Python Stack}
  \begin{columns}[T]
    \begin{column}{0.5\textwidth}
      \footnotesize
      \begin{tabularx}{\textwidth}{lX}
        \toprule
        \texttt{numpy} & arrays, linear algebra \\
        \texttt{pandas} & tables (\texttt{DataFrame}), Excel/CSV I/O \\
        \texttt{matplotlib} & plotting, PDF/SVG export \\
        \texttt{scipy} & statistics, optimisation \\
        \texttt{jupyterlab} & notebooks in the browser \\
        \midrule
        \texttt{apsg} & orientation data, tensors \\
        \texttt{petropandas} & mineral chemistry \\
        \bottomrule
      \end{tabularx}
    \end{column}
    \begin{column}{0.46\textwidth}
      \begin{block}{Why bother?}
        \small
        \begin{itemize}
          \item Every step is written down and re-runnable
          \item Raw data are never edited by hand
          \item 10 or 10\,000 measurements --- same code
          \item Figures regenerate when data change
          \item Free, open, and shared with the whole community
        \end{itemize}
      \end{block}
    \end{column}
  \end{columns}
\end{frame}

\section{Installation}

\begin{frame}[fragile]{Install with \texttt{uv}}
  \small
  \textbf{1. Install uv} (one-time; a fast Python project manager):
  \begin{lstlisting}[style=shell]
# Linux / macOS
curl -LsSf https://astral.sh/uv/install.sh | sh
# Windows (PowerShell)
powershell -ExecutionPolicy ByPass -c "irm https://astral.sh/uv/install.ps1 | iex"
  \end{lstlisting}
  \textbf{2a. Use the workshop materials} (recommended --- exact tested versions):
  \begin{lstlisting}[style=shell]
cd UPSG2026_workshop    # cloned or unzipped (...-main) materials
uv sync
  \end{lstlisting}
  \textbf{2b. \ldots or start your own project from scratch:}
  \begin{lstlisting}[style=shell]
uv init --bare --python 3.14 myproject && cd myproject
uv add "apsg[lab]" "petropandas[lab]"
  \end{lstlisting}
  \textbf{3. Start Jupyter Lab:} \texttt{uv run jupyter lab}
\end{frame}

\begin{frame}[fragile]{Check the Installation}
  Open a new notebook and run:
  \begin{lstlisting}[style=python]
import apsg, petropandas
print(apsg.__version__, petropandas.__version__)   # 2.0.5 0.2.3

from apsg import *
quicknet(fol(120, 40), lin(210, 30))                # a stereonet should appear

from petropandas import pd, Grt
from petropandas.data import minerals
minerals.oxides.select("Garnet", on="Mineral").mineral.end_members(Grt).head()
  \end{lstlisting}
  \vspace{0.2cm}
  \small
  \begin{itemize}
    \item Notebooks live in \texttt{notebooks/}, example data in \texttt{notebooks/data/},
          your results go to \texttt{notebooks/output/}
    \item Problems? Raise your hand now --- everything later builds on this
  \end{itemize}
\end{frame}

\section{Two Philosophies}

\begin{frame}[fragile]{\texttt{apsg}: Object-Oriented}
  \begin{columns}[T]
    \begin{column}{0.55\textwidth}
      \begin{lstlisting}[style=python]
from apsg import *

f = fol(250, 40)           # a plane
l = f ** fol(160, 30)      # intersection line
S = folset.random_fisher(
      position=f, kappa=20)  # a set of planes
S.ortensor().eigenlins()   # principal axes

F = defgrad.from_comp(xx=2, zz=0.5)
S.transform(F)             # deform the set

s = StereoNet()
s.great_circle(S); s.line(l); s.show()
      \end{lstlisting}
    \end{column}
    \begin{column}{0.42\textwidth}
      \small
      \begin{itemize}
        \item Data are \textbf{objects that know what they are}: a
              \texttt{lin} is not a \texttt{fol}, a \texttt{fault} has a sense
        \item Operators have geological meaning: \texttt{**} intersection,
              \texttt{@} tensor product
        \item Sets carry statistics; tensors transform anything
        \item Stereonets understand every object type
      \end{itemize}
    \end{column}
  \end{columns}
\end{frame}

\begin{frame}[fragile]{\texttt{petropandas}: DataFrame Integration}
  \begin{columns}[T]
    \begin{column}{0.55\textwidth}
      \begin{lstlisting}[style=python]
from petropandas import pd, Grt

df = pd.read_excel("epma_export.xlsx")
df.oxides()                # oxide columns
df.moles()                 # molar proportions
df.cations(n_oxygens=12)   # APFU
df.mineral.end_members(Grt)
df.mineral.check_stoichiometry(Grt)
      \end{lstlisting}
    \end{column}
    \begin{column}{0.42\textwidth}
      \small
      \begin{itemize}
        \item No new data structure --- your table \textbf{stays a
              \texttt{DataFrame}}
        \item Functionality hangs off \textbf{accessors}:
              \texttt{df.oxides}, \texttt{df.moles}, \texttt{df.cations},
              \texttt{df.mineral}, \texttt{df.bulk}
        \item All of pandas still works: filtering, grouping, \texttt{to\_excel}
      \end{itemize}
    \end{column}
  \end{columns}
  \vspace{0.3cm}
  \begin{block}{Bridge}
    \texttt{apsg} also has DataFrame accessors: \texttt{df.fol()}, \texttt{df.lin()},
    \texttt{df.fault()} turn table columns into \texttt{apsg} sets.
  \end{block}
\end{frame}

\begin{frame}{Materials}
  \centering
  \small
  \begin{tabularx}{0.92\textwidth}{lX}
    \toprule
    \textbf{Notebook} & \textbf{Content} \\
    \midrule
    \texttt{01\_apsg\_basics} & Block 1: features, stereonets, statistics, fold axes \\
    \texttt{02\_apsg\_advanced} & Block 2: faults, strain, progressive deformation, stress inversion \\
    \texttt{03\_petropandas} & Block 3: EPMA export $\rightarrow$ formulas, end-members, QC, charts \\
    \texttt{04a\_byod\_structural} & Block 4 template: your field data $\rightarrow$ stereonets + statistics \\
    \texttt{04b\_byod\_epma} & Block 4 template: your microprobe data $\rightarrow$ tables + charts \\
    \bottomrule
  \end{tabularx}

  \vspace{0.4cm}
  \raggedright
  Each teaching notebook ends with a hands-on exercise and a worked solution.
  Example data: \texttt{notebooks/data/}.
\end{frame}

\end{document}
